\documentclass[11pt]{article}
\usepackage[left=0.7in,right=0.7in,top=1in,bottom=0.7in]{geometry}
\usepackage{fancyhdr} % for header
\usepackage{graphicx} % for figures
\usepackage{amsmath}  % for extended math markup
\usepackage{amssymb}  % "" (more math symbols)
\usepackage[style=authoryear]{biblatex}
\usepackage[bookmarks=false]{hyperref} % for URL embedding
\usepackage[noend]{algpseudocode} % for pseudocode
\usepackage[plain]{algorithm} % float environment for algorithms
\usepackage{float}
\usepackage{mhchem}
\usepackage{xurl}
\newcommand{\bvec}[1]{\textbf{#1}}

\title{Assignment 2 Solutions}
\author{Charlie Li}
\date{\today}

\begin{document}

\maketitle

\section{Understanding word2vec}
\subsection*{(a)}
$y_w$ is 0 everywhere except for $y_o$, which equals 1.

\subsection*{(b)}
\begin{enumerate}
  \item 
  \begin{align*}
  \frac{\partial}{\partial \bvec{v}_c} \bvec{J}_{\text{naive-softmax}}(\bvec{v}_c, o, \bvec{U})
    &= \frac{\partial}{\partial \bvec{v}_c} \left(-\log P(O=o \mid C=c)\right) \\
    &= - \frac{\partial}{\partial \bvec{v}_c} \log \frac{\exp(\bvec{u}_o^{\top}\bvec{v}_c)}{\sum_{w\in \text{Vocab}} \exp(\bvec{u}_w^{\top}\bvec{v}_c)} \\
    &=  -\frac{\sum_{w\in \text{Vocab}} \exp(\bvec{u}_w^{\top}\bvec{v}_c)}{\exp(\bvec{u}_o^{\top}\bvec{v}_c)}\times \frac{\partial}{\partial \bvec{v}_c} \frac{\exp(\bvec{u}_o^{\top}\bvec{v}_c)}{\sum_{w\in \text{Vocab}} \exp(\bvec{u}_w^{\top}\bvec{v}_c)} \\
  \frac{\partial}{\partial \bvec{v}_c} \frac{\exp(\bvec{u}_o^{\top}\bvec{v}_c)}{\sum_{w\in \text{Vocab}} \exp(\bvec{u}_w^{\top}\bvec{v}_c)}  
    &= \frac{
      \frac{\partial}{\partial \bvec{v}_c} \exp(\bvec{u}_o^{\top}\bvec{v}_c) \sum_{w\in \text{Vocab}} \exp(\bvec{u}_w^{\top}\bvec{v}_c)
      - \exp(\bvec{u}_o^{\top}\bvec{v}_c) \frac{\partial}{\partial \bvec{v}_c} \sum_{w\in \text{Vocab}} \exp(\bvec{u}_w^{\top}\bvec{v}_c)
    }{
      \left(\sum_{w\in \text{Vocab}} \exp(\bvec{u}_w^{\top}\bvec{v}_c)\right)^2
    } \\
    &= \frac{
      \bvec{u}_o^{\top}\exp(\bvec{u}_o^{\top}\bvec{v}_c) \sum_{w\in \text{Vocab}} \exp(\bvec{u}_w^{\top}\bvec{v}_c) - \exp(\bvec{u}_o^{\top}\bvec{v}_c) \sum_{w\in \text{Vocab}} \bvec{u}_w^{\top} \exp(\bvec{u}_w^{\top}\bvec{v}_c)
    }{
      \left(\sum_{w\in \text{Vocab}} \exp(\bvec{u}_w^{\top}\bvec{v}_c)\right)^2
    }\\
    &= \frac{\bvec{u}_o^{\top}\exp(\bvec{u}_o^{\top}\bvec{v}_c)}{\sum_{w\in \text{Vocab}} \exp(\bvec{u}_w^{\top}\bvec{v}_c)} - \frac{ \exp(\bvec{u}_o^{\top}\bvec{v}_c) \sum_{w\in \text{Vocab}} \bvec{u}_w^{\top} \exp(\bvec{u}_w^{\top}\bvec{v}_c)}{\left(\sum_{w\in \text{Vocab}} \exp(\bvec{u}_w^{\top}\bvec{v}_c)\right)^2}\\
  \frac{\partial}{\partial \bvec{v}_c} \bvec{J}_{\text{naive-softmax}}(\bvec{v}_c, o, \bvec{U})
    &= -\bvec{u}_o^{\top} + \frac{\sum_{w\in \text{Vocab}} \bvec{u}_w^{\top} \exp(\bvec{u}_w^{\top}\bvec{v}_c)}{\sum_{w\in \text{Vocab}} \exp(\bvec{u}_w^{\top}\bvec{v}_c)}\\
    &= -\bvec{u}_o^{\top} + E_o[\bvec{u}_o^{\top}|C=c]
  \end{align*}
  With shape convention, \[ \frac{\partial}{\partial \bvec{v}_c} \bvec{J}_{\text{naive-softmax}}(\bvec{v}_c, o, \bvec{U}) = -\bvec{u}_o + E_o[\bvec{u}_o|C=c]   = -\bvec{U}y +\bvec{U}\hat{y} = \bvec{U}(-y+\hat{y}) \]

  \item 
  \[\bvec{U}(-y+\hat{y}) = 0 \]
  \[ -y+\hat{y}\in \text{Null}(\bvec{U}) \]
  
  \item Adding $\bvec{u}_o$ to the word vector $v_c$ makes it more aligned to $\bvec{u}_o$, which increases the dot product between $\bvec{u}_o$ and $\bvec{v}_c$, thus increasing the probability $P(O=o|C=c)$. Subtracting the expectation of $\bvec{u}_o$ zeros out the current guess and pushes the vector away from wrong outside words.
\end{enumerate}

\subsection*{(c)} It would affect the classification when a word that has a stronger emotion also has a word vector with a larger magnitude. It would not affect the classification when the length of the word vector is not correlated with the strength of the meaning of the word.

\subsection*{(d)}
$$\frac{\partial J}{\partial U} = \left[ \frac{\partial J}{\partial \bvec{u}_1}, \frac{\partial J}{\partial \bvec{u}_2}, \dots, \frac{\partial J}{\partial \bvec{u}_{|\text{Vocab}|}} \right]$$

\subsection*{(e)}
\begin{align*}
  \frac{\partial J}{\partial \bvec{u}_{i}} &= -\frac{\partial}{\partial \bvec{u}_{i}} log\frac{exp(\bvec{u}_o^{\top}\bvec{v}_c)}{\sum_{w\in \text{Vocab}}exp(\bvec{u}_w^{\top}\bvec{v}_c)}\\
    &= -\frac{\sum_{w\in \text{Vocab}}exp(\bvec{u}_w^{\top}\bvec{v}_c)}{exp(\bvec{u}_o^{\top}\bvec{v}_c)}\\
    & \left[ \frac{
      \frac{\partial}{\partial \bvec{u}_{i}}exp(\bvec{u}_o^{\top}\bvec{v}_c)\sum_{w\in \text{Vocab}}exp(\bvec{u}_w^{\top}\bvec{v}_c) - exp(\bvec{u}_o^{\top}\bvec{v}_c)\frac{\partial}{\partial \bvec{u}_{i}}(\sum_{w\in \text{Vocab}}exp(\bvec{u}_w^{\top}\bvec{v}_c))
    }{(\sum_{w\in \text{Vocab}}exp(\bvec{u}_w^{\top}\bvec{v}_c))^2} \right]\\
    &=AB \\
  \text{case 1: }i\ne o\\
  B &= \frac{-exp(\bvec{u}_o^{\top}\bvec{v}_c)\bvec{v}_c exp(\bvec{u}_i^{\top}\bvec{v}_c)}{(\sum_{w\in \text{Vocab}}exp(\bvec{u}_w^{\top}\bvec{v}_c))^2}\\
  AB &= \frac{\bvec{v}_c exp(\bvec{u}_i^{\top}\bvec{v}_c)}{\sum_{w\in \text{Vocab}}exp(\bvec{u}_w^{\top}\bvec{v}_c)}\\
    &= \bvec{v}_c \hat{y}_i\\
  \text{case 2: }i = o\\
  B &= \frac{\bvec{v}_cexp(\bvec{u}_o^{\top}\bvec{v}_c)}{\sum_{w\in \text{Vocab}}exp(\bvec{u}_w^{\top}\bvec{v}_c)} + B_{\text{case 1}}\\
  AB &= -\bvec{v}_c + \bvec{v}_c \hat{y}_i\\
  \text{unified equation}\\
  \frac{\partial J}{\partial \bvec{u}_{i}} &= \bvec{v}_c(\hat{y}_i - y_i)
\end{align*}

\subsection*{(f)}
\begin{equation*}
  f'(x) = \begin{cases}
    1 & x>0\\
    \alpha & x<0
  \end{cases}
\end{equation*}

\subsection*{(g)}
\begin{align*}
  \sigma'(x) &= \frac{(e^x)'(e^x+1)-e^x(e^x+1)'}{(e^x+1)^2} \\
  &= \frac{e^x(e^x+1) - e^x e^x}{(e^x+1)^2}\\
  &= \frac{e^x}{e^{2x}+1+2e^x}\\
  &= \frac{1}{e^x+e^{-x}+2}\\
  1+e^{-x} &= \frac{1}{\sigma(x)}\\
  e^{-x} &= \frac{1-\sigma(x)}{\sigma(x)}\\
  e^{x} &= \frac{\sigma(x)}{1-\sigma(x)}\\
  \sigma'(x) &= \frac{1}{\frac{\sigma(x)}{1-\sigma(x)}+\frac{1-\sigma(x)}{\sigma(x)}+2}\\
  &= \frac{\sigma(x)(1-\sigma(x))}{\sigma^2(x)+(1-\sigma(x))^2+2\sigma(x)(1-\sigma(x))}\\
  &= \sigma(x)(1-\sigma(x))
\end{align*}

\section{Neural Networks Optimization}

\subsection*{(a)}
\begin{enumerate}
  \item The new gradient is multiplied by a coefficient $1-\beta_1$, which makes the change smaller. It accumulates the history of the gradients, which prevents the model from being trapped in shallow local minima.
  \item Parameters with smaller $v$ will get larger updates, which corresponds to smaller gradient of $J$ with respect to that parameter. This helps with learning by helping parameters get out of a plateau or reach a minimum faster.
\end{enumerate}

\subsection*{(b)}
\begin{enumerate}
  \item \begin{align*}
    h_{drop, i} &= \lambda d h_i\\
    E[h_{drop, i}] &= E[\lambda d h_i] \\
      &= \lambda E[d] h_i \\
      &= h_i \\
    \lambda &= \frac{1}{E[d]} \\
      &= \frac{1}{p_{drop}}
  \end{align*}

\section{Neural Transition-Based Dependency Parsing}
\subsection*{(a)}
\begin{table}[htbp]
  \centering
  \small
  \renewcommand{\arraystretch}{1.3}
  \resizebox{\textwidth}{!}{
    \begin{tabular}{|l|l|l|l|}
      \hline
      \textbf{Stack} & \textbf{Buffer} & \textbf{New dependency} & \textbf{Transition} \\
      \hline
      [ROOT, presented, my] & [findings, at, the, NLP, conference] & & \texttt{SHIFT}\\
      \hline
      [ROOT, presented, my, findings] & [at, the, NLP, conference] & & \texttt{SHIFT}\\
      \hline
      [ROOT, presented, findings] & [at, the, NLP, conference] & findings $\rightarrow$ my & \texttt{LEFT-ARC}\\
      \hline
      [ROOT, presented] & [at, the, NLP, conference] & presented $\rightarrow$ findings & \texttt{RIGHT-ARC}\\
      \hline
      [ROOT, presented, at] & [the, NLP, conference] & & \texttt{SHIFT}\\
      \hline
      [ROOT, presented, at, the] & [NLP, conference] & & \texttt{SHIFT}\\
      \hline
      [ROOT, presented, at, the, NLP] & [conference] & & \texttt{SHIFT}\\
      \hline
      [ROOT, presented, at, the, NLP, conference] & [] & & \texttt{SHIFT}\\
      \hline
      [ROOT, presented, at, the, conference] & [] & conference $\rightarrow$ NLP & \texttt{LEFT-ARC}\\
      \hline
      [ROOT, presented, at, conference] & [] & conference $\rightarrow$ the & \texttt{LEFT-ARC}\\
      \hline
      [ROOT, presented, conference] & [] & conference $\rightarrow$ at & \texttt{LEFT-ARC}\\
      \hline
      [ROOT, presented] & [] & presented $\rightarrow$ conference & \texttt{RIGHT-ARC}\\
      \hline
      [ROOT] & [] & ROOT $\rightarrow$ presented & \texttt{RIGHT-ARC}\\
      \hline
    \end{tabular}
  }
\end{table}

\subsection*{(b)}
$2n$ steps. It takes $n$ shifts to empty the buffer, and $n$ right or left arcs to empty the stack to one element.

\subsection*{(e)}
\subsubsection*{i}
  \begin{align*}
    h_i &= ReLU(\bvec{x}\bvec{W}_i+b_{1,i}) \text{  where $\bvec{W}_i$ is column $i$ in matrix $\bvec{W}$}\\
    \frac{\partial h_i}{\partial x_j} &= \frac{\partial ReLU(\bvec{x}\bvec{W}_i+b_{1,i})}{\partial \bvec{x}\bvec{W}_i+b_{1,i}} \frac{\partial \bvec{x}\bvec{W}_i+b_{1,i}}{\partial x_j}\\
      &= [(\bvec{x}\bvec{W}_i+b_{1,i}) > 0]W_{j, i}
  \end{align*}

\subsubsection*{ii}
  \begin{align*}
    \frac{\partial J}{\partial \hat{y}_i} &= -\frac{y_i}{\hat{y}_i}\\
    \hat{y}_i &= \frac{e^{l_i}}{\sum_{j=1}^{3}e^{l_j}}\\
    \frac{\partial \hat{y}_i}{\partial l_i} &= \frac{e^{l_i}(\sum_{j=1}^{3}e^{l_j}) - e^{l_i}e^{l_i}}{(\sum_{j=1}^{3}e^{l_j})^2}\\
    \frac{\partial J}{\partial \hat{l}_i} &=  -\frac{y_i}{\hat{y}_i}\frac{e^{l_i}(\sum_{j=1}^{3}e^{l_j}) - e^{2l_i}}{(\sum_{j=1}^{3}e^{l_j})^2}\\
  \end{align*}

\subsubsection*{iii}
\begin{itemize}
  \item UAS on dev set: 72.10
  \item UAS on test set: 89.31
\end{itemize}

\end{enumerate}

\end{document}
